\chapter{Event-Driven Backtests}\label{ch:rs:event-driven-backtests}

A strategy that rests limit orders one tick either side of the last price fills, in a \omterm{def:rs:market-data-for-research:bar}{bar} \omterm{def:rs:vectorised-backtests:backtest}{backtest}, every time a \omterm{def:rs:market-data-for-research:bar}{bar}'s range reaches its
price. Asked instead to see the price trade one tick through its order before counting a fill, the same \omterm{def:rs:vectorised-backtests:backtest}{backtest} fills 55\% as often and turns
a small daily profit into a loss of \$584 a day, with fills followed by a two-tick move against it. Add an order-entry latency of five seconds,
and the \omterm{def:rs:vectorised-backtests:backtest}{backtest}'s answer depends on a question the \omterm{def:rs:market-data-for-research:bar}{bars} cannot answer: did the price touch the order while it was live? Level~1 of
chapter~16 cannot even express such a strategy. The \omterm{def:rs:event-driven-backtests:evbt}{event-driven backtest} can: it simulates orders on a clock, and it makes every assumption
about fills an explicit, replaceable model. This chapter builds it (\texttt{firm.evbt}) and shows how much of a passive strategy's result is the
\omterm{def:rs:event-driven-backtests:fill}{fill model}'s.

\section{Architecture: clock, events, handlers}

\begin{definition}[Event-driven backtest, simulated clock]\label{def:rs:event-driven-backtests:evbt}
An \emph{event-driven backtest}\index{event-driven backtest} processes a time-ordered stream of events (market data, order acknowledgements,
fills, cancellations, corporate actions) one at a time, each changing the state of the strategy, its orders and its portfolio. Its
\emph{simulated clock}\index{simulated clock} is the time of the event being processed; nothing may use information stamped later.
\end{definition}

\begin{omfigure}
\begin{tikzpicture}[font=\small,
  box/.style={draw=omDef,fill=omDef!8,rounded corners=2pt,align=center,inner sep=4pt,minimum height=9mm,text width=2.7cm},
  arr/.style={-Stealth,thick,draw=black!60}]
\node[box] (q) at (0,0) {event queue\\ (time, kind, sequence)};
\node[box] (d) at (3.6,0) {dispatcher:\\ advances the clock};
\node[box] (m) at (7.2,0) {order manager:\\ lifecycle};
\node[box] (p) at (10.8,0) {portfolio:\\ position, cash};
\node[box] (s) at (7.2,2.0) {strategy:\\ on\_bar, on\_fill};
\node[box] (f) at (7.2,-2.0) {fill model: touch,\\ penetration, cap};
\draw[arr] (q) -- (d);
\draw[arr] (d) -- (m);
\draw[arr] (m) -- (p);
\draw[arr] (d) |- (s);
\draw[arr] (s) -- node[right,font=\footnotesize]{submit, cancel} (m);
\draw[arr] (m) -- (f);
\draw[arr] (f) -| (p);
\draw[arr,dashed] (s.north) -- ++(0,0.5) -| node[pos=0.25,above,font=\footnotesize]{acknowledgements and cancellations, after the latency} (q.north);
\end{tikzpicture}

\omcaption{The architecture of \texttt{firm.evbt}: one queue of events ordered by time, kind and sequence number, a dispatcher, and handlers that
may only schedule events in the future.}\label{fig:rs:event-driven-backtests:architecture}
\end{omfigure}

\texttt{firm.evbt} keeps one priority queue. At each \omterm{def:rs:market-data-for-research:bar}{bar}'s end the dispatcher first matches the working orders against the \omterm{def:rs:market-data-for-research:bar}{bar} through the \omterm{def:rs:event-driven-backtests:fill}{fill
model}, then marks the portfolio, then calls the strategy, which may submit or cancel orders; a submission becomes a working order only when its
acknowledgement event arrives, after the order-entry latency, and a cancellation takes effect only when its own event does. Events at the same
time are ordered by kind (\omterm{def:rs:market-data-for-research:bar}{bars} before order events) and then by a sequence number, so a run is deterministic: the same data, strategy and
models give the same fills to the last digit. The result is the \texttt{BacktestResult} of chapter~16, with the orders and fills attached.

\section{Orders and fills at bar level}

\begin{definition}[Order lifecycle, partial fill]\label{def:rs:event-driven-backtests:lifecycle}
The \emph{order lifecycle}\index{order lifecycle} is the sequence of states an order passes through: pending (sent, not yet acknowledged),
working, partially filled, and one of filled or cancelled. A \emph{partial fill}\index{partial fill} executes part of an order's quantity, leaving
the rest working.
\end{definition}

\begin{definition}[Fill model, touch fill, penetration fill, volume participation cap]\label{def:rs:event-driven-backtests:fill}
A \emph{fill model}\index{fill model} decides from market data whether a working order executes, at what price and for what quantity. The
\emph{touch fill}\index{touch fill} fills a limit order in full when a \omterm{def:rs:market-data-for-research:bar}{bar} trades at its price; the \emph{penetration fill}\index{penetration fill}
only when the \omterm{def:rs:market-data-for-research:bar}{bar} trades a given number of ticks beyond it; a \emph{volume participation cap}\index{volume participation cap} limits the quantity
filled in a \omterm{def:rs:market-data-for-research:bar}{bar} to a share of the \omterm{def:rs:market-data-for-research:bar}{bar}'s volume.
\end{definition}

A \omterm{def:rs:market-data-for-research:bar}{bar} records four prices and a volume; it does not record where the strategy's order stood in the queue at its price, how much traded there, or
whether the trades at that price came before or after the order arrived. Each \omterm{def:rs:event-driven-backtests:fill}{fill model} fills that gap with an assumption. The \omterm{def:rs:event-driven-backtests:fill}{touch fill} assumes
the order was first in the queue. The \omterm{def:rs:event-driven-backtests:fill}{penetration fill} assumes it was last: if the price traded through, everything ahead of it must have traded.
The volume cap assumes the order can take a fixed share of the \omterm{def:rs:market-data-for-research:bar}{bar}'s volume; the open-source backtester Zipline's \texttt{VolumeShareSlippage}
model, for instance, caps fills at 2.5\% of each \omterm{def:rs:market-data-for-research:bar}{bar}'s volume by default and prices them off the \omterm{def:rs:market-data-for-research:bar}{bar}'s close with a quadratic impact. The truth lies
between the touch and \omterm{def:rs:event-driven-backtests:fill}{penetration fills} and depends on the queue (chapter~18).

The chapter's strategy (\texttt{rs\_evbt.Quoter}) runs on one-minute \omterm{def:rs:market-data-for-research:bar}{bars} built from four simulated 6.5-hour days of \texttt{firm.tape}: at every
\omterm{def:rs:market-data-for-research:bar}{bar}'s close it cancels its orders and rests one lot (100 shares) one tick below the close and one lot one tick above, within an inventory of five
lots, marking to the close. Per day, on average:

\begin{center}
\small
\begin{tabular}{@{}lcccc@{}}
\toprule
\omterm{def:rs:event-driven-backtests:fill}{fill model} (no latency) & lots filled & fill rate & P\&L (\$) & mark-out (ticks)\\
\midrule
touch & 523 & 70\% & $+17$ & $-0.16$\\
penetration by one tick & 290 & 41\% & $-584$ & $-2.00$\\
touch, capped at 2.5\% of volume & 388 & 55\% & $+20$ & $-0.18$\\
level 1: buy after a down \omterm{def:rs:market-data-for-research:bar}{bar}, sell after an up \omterm{def:rs:market-data-for-research:bar}{bar} & --- & --- & $-215$ & ---\\
\bottomrule
\end{tabular}
\end{center}

The mark-out (the close five \omterm{def:rs:market-data-for-research:bar}{bars} after a fill, in the fill's direction) measures adverse selection: the \omterm{def:rs:event-driven-backtests:fill}{penetration fills} are the ones that happen
when the price is moving through the order, and they lose two ticks on average in five minutes (\cref{fig:rs:event-driven-backtests:models}). The \omterm{def:rs:event-driven-backtests:fill}{touch
fill}'s daily P\&L (from $-\$333$ to $+\$206$ over the four days) is within its noise of zero; the \omterm{def:rs:event-driven-backtests:fill}{penetration fill}'s is negative every day. A fill rate
of 70\% for an order one tick away from the price is itself a warning: a real order joins the back of a queue that may hold thousands of shares.

\section{Decision time and latency}

\begin{definition}[Latency and the live-in-bar policy]\label{def:rs:event-driven-backtests:latency}
With an order-entry latency, an order sent at a \omterm{def:rs:market-data-for-research:bar}{bar}'s close is live from some moment within the next \omterm{def:rs:market-data-for-research:bar}{bar}. A \emph{conservative} policy lets an order
fill in a \omterm{def:rs:market-data-for-research:bar}{bar} only if it was live for the whole \omterm{def:rs:market-data-for-research:bar}{bar}; an \emph{optimistic} policy, if it was live at any moment of it. A \omterm{def:rs:market-data-for-research:bar}{bar} cannot say which is right,
because it does not say when within it the price reached the order.
\end{definition}

With a latency of five seconds, the quoter's orders, cancelled and replaced at every \omterm{def:rs:market-data-for-research:bar}{bar}'s close, are never live for a whole \omterm{def:rs:market-data-for-research:bar}{bar}: under the
conservative policy they never fill. Under the optimistic policy they fill more often than with no latency (585 lots a day for the \omterm{def:rs:event-driven-backtests:fill}{touch fill}), and
lose more: $-\$41$ a day for the \omterm{def:rs:event-driven-backtests:fill}{touch fill}, $-\$776$ for the \omterm{def:rs:event-driven-backtests:fill}{penetration fill}, $-\$190$ for the capped fill, with the capped fills' mark-out falling
from $-0.18$ to $-0.87$ ticks. The two policies bracket the answer, and a strategy whose result swings from nothing to a loss between them cannot be
judged at level~2. That is the signal to move to level~3, where the order's arrival, its queue position and every trade at its price are events of their
own.

\begin{omfigure}
\begin{tikzpicture}
\begin{axis}[name=lots,width=7cm,height=5.4cm,ybar,bar width=9pt,title={\small lots filled a day},tick label style={font=\footnotesize},
  xtick={0,1,2,3,4,5},xticklabels={touch,penetr.,capped,touch,penetr.,capped},x tick label style={rotate=45,anchor=east},
  xmin=-0.6,xmax=5.6,ymin=0,ymax=760,grid=major,grid style={gray!25},table/col sep=comma]
\addplot[fill=omDef!55,draw=omDef] table[x=k,y=lots]{figdata/research/17-event-driven-backtests/models.csv};
\draw[black!50,dashed] (axis cs:2.5,0) -- (axis cs:2.5,760);
\node[font=\footnotesize,anchor=north] at (axis cs:1,750) {no latency};
\node[font=\footnotesize,anchor=north] at (axis cs:4,750) {5 s, optimistic};
\end{axis}
\begin{axis}[at={(lots.east)},anchor=west,xshift=1.4cm,width=7cm,height=5.4cm,ybar,bar width=9pt,title={\small P\&L a day (\$)},
  tick label style={font=\footnotesize},xtick={0,1,2,3,4,5,6},xticklabels={touch,penetr.,capped,touch,penetr.,capped,level 1},
  x tick label style={rotate=45,anchor=east},ymin=-850,ymax=100,grid=major,grid style={gray!25},table/col sep=comma]
\addplot[fill=omThm!45,draw=omThm] table[x=k,y=pnl]{figdata/research/17-event-driven-backtests/models.csv};
\draw[black!50,dashed] (axis cs:2.5,-850) -- (axis cs:2.5,100);
\end{axis}
\end{tikzpicture}

\omcaption{The passive quoter on four simulated days under each \omterm{def:rs:event-driven-backtests:fill}{fill model}, with no latency and with a five-second latency under the optimistic
policy (under the conservative policy it never fills), and the level-1 version of the same idea. Data: \texttt{rs\_evbt} on \texttt{firm.tape} \omterm{def:rs:market-data-for-research:bar}{bars}.}\label{fig:rs:event-driven-backtests:models}
\end{omfigure}

\section{Calendars and corporate actions in the loop}

Events that a \omterm{def:rs:vectorised-backtests:backtest}{vectorised backtest} folds into adjusted prices are, in an event-driven one, events in the queue. A split multiplies the position and the
working orders' quantities by its ratio and divides their limit prices by it, at the moment it takes effect; the \omterm{def:rs:market-data-for-research:bar}{bars} keep the prices as traded, so a
limit order that survived a split unadjusted would sit at twice the new price and fill instantly. A dividend pays cash on its payment date to positions
held at the record date. Half-days end the session early, auctions replace continuous trading at the open and the close (Book~1, chapter~13), and a halt
freezes the book; each belongs in the calendar the dispatcher reads, with its own handler, and each has caused a \omterm{def:rs:vectorised-backtests:backtest}{backtest} to trade at a time the market
was closed. \texttt{firm.evbt} implements splits and leaves the others as the build's stretch; its acceptance test holds a lot through a two-for-one
split and checks that the position doubles and the capital does not move.

\section{Determinism}

A \omterm{def:rs:vectorised-backtests:backtest}{backtest} that gives different answers on two runs cannot be debugged, compared or reviewed. Determinism needs three things: a total order on events
(time, then kind, then sequence number; never the order in which a dictionary happens to iterate), no randomness except from seeded generators passed
in explicitly, and no dependence on the wall clock. \texttt{firm.evbt}'s tests run the same \omterm{def:rs:vectorised-backtests:backtest}{backtest} twice and compare the net returns bit for bit. The
same discipline lets a level-2 and a level-3 \omterm{def:rs:vectorised-backtests:backtest}{backtest} be compared on the same events, and a live run be replayed against its own logs (chapter~19).

\section{Tutorial: the limit order that always filled}\label{tut:rs:event-driven-backtests:tut}

\begin{tutorial}
\textbf{Goal.} Run the passive quoter in \texttt{firm.evbt} under three \omterm{def:rs:event-driven-backtests:fill}{fill models} and two latencies with both live-in-bar policies, and compare with
the level-1 version of the same idea. \textbf{End state:} \cref{fig:rs:event-driven-backtests:models}; the table.
\begin{tutsteps}
\item \textbf{The \omterm{def:rs:event-driven-backtests:fill}{fill models}}: touch, penetration, and a cap on the \omterm{def:rs:market-data-for-research:bar}{bar}'s volume.
\omcode{code/firm/evbt/firm_evbt.py}{70}{104}{Three fill models for bars.}\label{lst:rs:event-driven-backtests:fills}
\item \textbf{Matching}: which orders were live in the \omterm{def:rs:market-data-for-research:bar}{bar}, and what the model fills.
\omcode{code/firm/evbt/firm_evbt.py}{151}{175}{Matching working orders against a bar.}\label{lst:rs:event-driven-backtests:match}
\item \textbf{Run} \texttt{summary()}, \texttt{level1()} and \texttt{fig\_evbt.py}.
\end{tutsteps}
\textbf{What to change next.} Rest the orders at the close instead of one tick away and watch the touch model's fill rate approach 100\%; add a
cancel-only-if-the-price-moved rule so orders survive whole \omterm{def:rs:market-data-for-research:bar}{bars} under latency.
\end{tutorial}

\section{Build: the level-2 backtester}\label{bld:rs:event-driven-backtests:evbt}

\begin{build}
\bfield{Purpose} The firm's backtester for strategies whose orders, fills and timing matter at \omterm{def:rs:market-data-for-research:bar}{bar} resolution: limit orders, stop orders, intraday
schedules, corporate actions in the loop.
\bfield{Interface} \texttt{Order}, \texttt{Fill}, \texttt{TouchFill()}, \texttt{PenetrationFill(ticks)}, \texttt{VolumeCapFill(inner, participation)},
\texttt{Strategy.on\_bar(ctx, i, bar)}, \texttt{Strategy.on\_fill(ctx, fill)}, \texttt{Engine(bars, strategy, fill\_model, latency, fee, capital, splits,
policy).run()} returning the \texttt{BacktestResult}, the orders and the fills.
\bfield{Rules} Handlers schedule only future events; orders fill only after their acknowledgement; every fill assumption lives in a \omterm{def:rs:event-driven-backtests:fill}{fill model}; the
live-in-bar policy is stated; the run is deterministic.
\bfield{Acceptance tests} \texttt{code/firm/evbt/tests/}: touch, penetration and capped fills on hand \omterm{def:rs:market-data-for-research:bar}{bars} with their lifecycles; a market order at the next
open; a latency that misses a \omterm{def:rs:market-data-for-research:bar}{bar}; cancellations; the optimistic policy filling an order live for an instant; fees, P\&L and a split; two runs identical.
\bfield{Stretch} Dividends, half-days, auctions and halts as calendar events; stop orders; several instruments on one clock.
\end{build}

\begin{omsources}
\item Zipline (open-source Python backtester), \texttt{zipline/finance/slippage.py}: \texttt{VolumeShareSlippage}, default volume limit of 2.5\% of each bar.
\end{omsources}

\section{Exercises}

\begin{exercise}[$\star$]\label{exo:rs:event-driven-backtests:1}
A buy limit at 99.00 is working. The next \omterm{def:rs:market-data-for-research:bar}{bar} has open 99.20, high 99.30, low 99.00 and close 99.10. Does it fill under the \omterm{def:rs:event-driven-backtests:fill}{touch fill}? Under a one-tick
\omterm{def:rs:event-driven-backtests:fill}{penetration fill} with a tick of 0.01?
\end{exercise}

\begin{exercise}[$\star$]\label{exo:rs:event-driven-backtests:2}
An order for 5\,000 shares faces \omterm{def:rs:market-data-for-research:bar}{bars} of 40\,000, 10\,000 and 60\,000 shares with a 2.5\% volume cap. How much fills in each \omterm{def:rs:market-data-for-research:bar}{bar}?
\end{exercise}

\begin{exercise}[$\star$]\label{exo:rs:event-driven-backtests:3}
Draw the lifecycle of an order for 300 shares that fills 100, then is cancelled with 200 remaining.
\end{exercise}

\begin{exercise}[$\star\star$]\label{exo:rs:event-driven-backtests:4}
Why are \omterm{def:rs:event-driven-backtests:fill}{penetration fills} followed by adverse price moves more than \omterm{def:rs:event-driven-backtests:fill}{touch fills}? Relate it to the queue at the order's price.
\end{exercise}

\begin{exercise}[$\star\star$]\label{exo:rs:event-driven-backtests:5}
With a latency of 5 seconds and one-minute \omterm{def:rs:market-data-for-research:bar}{bars}, why does the quoter never fill under the conservative policy? How would a strategy that cancels only when
the price moves fare?
\end{exercise}

\begin{exercise}[$\star\star$]\label{exo:rs:event-driven-backtests:6}
A two-for-one split takes effect while a sell limit at \$120 for 100 shares is working and the position is 100 shares. What must the engine do, and what
goes wrong if it does not?
\end{exercise}

\begin{exercise}[$\star\star\star$]\label{exo:rs:event-driven-backtests:7}
\emph{Coding.} Run the quoter with the orders at the close price (not one tick away) under the touch and \omterm{def:rs:event-driven-backtests:fill}{penetration fills}. What happens to the fill rates and
the mark-outs?
\end{exercise}

\begin{exercise}[$\star\star\star$]\label{exo:rs:event-driven-backtests:8}
\emph{Find the flaw.} ``Our \omterm{def:rs:vectorised-backtests:backtest}{backtest} fills limit orders when the \omterm{def:rs:market-data-for-research:bar}{bar}'s low reaches them, and the strategy made money every week for two years.''
\end{exercise}

\section{Problem: The Limit Order That Always Filled}

\begin{problem}[{Weekend problem --- a passive strategy, three fill models}]\label{pb:rs:event-driven-backtests:1}
The chapter's quoter on four simulated days of \texttt{firm.tape} one-minute \omterm{def:rs:market-data-for-research:bar}{bars}.

\medskip\noindent\textbf{Part I --- The engine.}
\begin{enumerate}
\item In what order does \texttt{firm.evbt} process a \omterm{def:rs:market-data-for-research:bar}{bar}'s close: fills, marks, strategy? Why that order?
\item When does a submitted order first become able to fill?
\item What makes two runs identical?
\item What does the engine return?
\end{enumerate}

\medskip\noindent\textbf{Part II --- \omterm{def:rs:event-driven-backtests:fill}{Fill models}.}
\begin{enumerate}[resume]
\item Give the lots filled, fill rate, P\&L and mark-out a day under each \omterm{def:rs:event-driven-backtests:fill}{fill model} with no latency.
\item Which assumption about the queue does each model make?
\item Why is the \omterm{def:rs:event-driven-backtests:fill}{penetration fill}'s P\&L negative every day while the \omterm{def:rs:event-driven-backtests:fill}{touch fill}'s is not?
\item What is the level-1 version of the idea, and what does it earn?
\end{enumerate}

\medskip\noindent\textbf{Part III --- Latency.}
\begin{enumerate}[resume]
\item What happens at a latency of 5 seconds under the conservative policy?
\item And under the optimistic policy, for each \omterm{def:rs:event-driven-backtests:fill}{fill model}?
\item What does the difference between the two policies tell you about the strategy?
\item How would you change the strategy so that latency matters less?
\end{enumerate}

\medskip\noindent\textbf{Part IV --- The verdict.}
\begin{enumerate}[resume]
\item State the \emph{named result}: the fill rate and P\&L of the passive \omterm{def:rs:market-data-for-research:bar}{bar} strategy under touch, penetration and capped fills, with and without
latency.
\item Which result would you report to a portfolio manager, and with which caveat?
\item What would a level-3 replay add (chapter~18)?
\item What does Zipline's default cap mean for a strategy that trades 5\% of volume?
\item Which corporate actions belong in the event queue for this strategy?
\item Why is a 70\% fill rate for orders one tick away suspicious?
\item What would you log from a live run to calibrate the \omterm{def:rs:event-driven-backtests:fill}{fill model}?
\item In one sentence: what is the \omterm{def:rs:event-driven-backtests:fill}{fill model}?
\end{enumerate}
\end{problem}

\section{Interview questions}

\begin{interviewq}[$\star$ \iqroles{developer, researcher}]\label{iq:rs:event-driven-backtests:1}
Describe the architecture of an event-driven backtester.
\end{interviewq}

\begin{interviewq}[$\star\star$ \iqroles{researcher, trader}]\label{iq:rs:event-driven-backtests:2}
How would you simulate limit-order fills on \omterm{def:rs:market-data-for-research:bar}{bar} data, and what biases does each approach have?
\end{interviewq}

\begin{interviewq}[$\star\star$ \iqroles{developer}]\label{iq:rs:event-driven-backtests:3}
How do you make a \omterm{def:rs:vectorised-backtests:backtest}{backtest} deterministic?
\end{interviewq}

\begin{interviewq}[$\star\star$ \iqroles{researcher}]\label{iq:rs:event-driven-backtests:4}
What is adverse selection for a passive order, and how would you measure it in a \omterm{def:rs:vectorised-backtests:backtest}{backtest}?
\end{interviewq}

\begin{interviewq}[$\star\star$ \iqroles{developer, trader}]\label{iq:rs:event-driven-backtests:5}
How should a backtester handle a stock split while orders are working?
\end{interviewq}

\begin{interviewq}[$\star\star\star$ \iqroles{researcher}]\label{iq:rs:event-driven-backtests:6}
When is a level-2 \omterm{def:rs:vectorised-backtests:backtest}{backtest} not enough? Give a test that tells you.
\end{interviewq}
